\subsubsection{The Linear-Algebra Interpretation}\label{sec:linear-algebra-interpretation-simons-algorithm}

In previous sections, we discussed the constraints that the hidden string must satisfy:
\begin{equation*}
    a \cdot y = 0 \pmod 2
\end{equation*}
Here we discuss the linear-algebra interpretation of this constraint. First of all, the condition is called \textbf{orthogonality} between the hidden string $a$ and the measured string $y$. We weren't introducing a new property of Simon's algorithm, but rather we were \textbf{renaming something we already knew}.

\highspace
\textcolor{Green3}{\faIcon{question-circle} \textbf{What does ``orthogonal'' mean here?}} We already know ordinary orthogonality. For real vectors, for example:
\begin{equation*}
    u = (1, 1), \quad v = (1, -1)
\end{equation*}
Their ordinary dot product is:
\begin{equation*}
    u \cdot v = 1 \cdot 1 + 1 \cdot (-1) = 0
\end{equation*}
So they are orthogonal. In Simon, however, our vectors are \textbf{bitstrings} and arithmetic is performed over:
\begin{equation*}
    GF(2) = \{0, 1\}
\end{equation*}
So for:
\begin{equation*}
    a = a_1 a_2 \ldots a_n, \quad y = y_1 y_2 \ldots y_n
\end{equation*}
We define the dot product as:
\begin{equation*}
    a \cdot y = a_1 y_1 \oplus a_2 y_2 \oplus \ldots \oplus a_n y_n
\end{equation*}
If the result is $a \cdot y = 0$, we say that \hl{$y$ is orthogonal to $a$ over $GF(2)$.}\footnote{%
    $GF(2)$ is the formal mathematical name for th arithmetic system we are using when we say ``modulo 2''. It is the acronym for \definition{Galois Field} with $2$ elements. The two elements, in our case, are simply $0$ and $1$.
} For example, if $a = 01$, for $y = 10$ we have:
\begin{equation*}
    a \cdot y = 01 \cdot 10 = 0(1) \oplus 1(0) = 0 \oplus 0 = 0
\end{equation*}
So $y = 10$ is orthogonal to $a = 01$:
\begin{equation*}
    01 \perp 10 \quad \text{over } GF(2)
\end{equation*}
This is \textbf{not geometric perpendicularity in ordinary Euclidean space}. It's algebraic orthogonality according to the mod-2 dot product.

\highspace
\textcolor{Green3}{\faIcon{question-circle} \textbf{If every measured $y$ must be orthogonal to $a$, which $y$'s are actually possible?}} That leads us to the \textbf{set of all possible Simon measurements}. We know that Simon can only produce $y$'s satisfying:
\begin{equation*}
    a \cdot y = 0
\end{equation*}
So for a fixed hidden string $a$, we define the set of all possible measurements as:
\begin{equation}
    a^{\perp} = \left\{
        y \in \left\{0,1\right\}^{n} \mid a \cdot y = 0
    \right\}
\end{equation}
In words, take \textbf{all possible $n$-bit strings} $y$ and keep only those satisfying the orthogonality condition. This is called the \textbf{orthogonal component} of $a$ \emph{over} $GF(2)$.

\highspace
For example, if $n=2$ and $a = 01$, there are four possible $y$'s:
\begin{equation*}
    00, 01, 10, 11
\end{equation*}
We can check which ones are orthogonal to $a$:
\begin{align*}
    01 \cdot 00 &= 0 \\
    01 \cdot 01 &= 1 \\
    01 \cdot 10 &= 0 \\
    01 \cdot 11 &= 1
\end{align*}
So the set of all possible measurements is:
\begin{equation*}
    a^{\perp} = \{00, 10\}
\end{equation*}
And that's exactly what we found explicitly running the circuit in the previous section. So the linear-algebra interpretation and the quantum interference calculation are describing \textbf{the same result}.

\highspace
\textcolor{Green3}{\faIcon{question-circle} \textbf{How large and how structured is the set of all possible measurements?}} We have just defined the ``\emph{orthogonal component}'' of $a$ over $GF(2)$, but \textbf{\emph{how many independent choices are there inside $a^{\perp}$?}} Let's look at the $n=3$ case. If $a = 111$, the condition for $y = y_1 y_2 y_3$ to be orthogonal to $a$ is:
\begin{equation*}
    a \cdot y = 0
\end{equation*}
Therefore, we have:
\begin{equation*}
    111 \cdot y_1 y_2 y_3 = 1(y_1) \oplus 1(y_2) \oplus 1(y_3) = y_1 \oplus y_2 \oplus y_3 = 0
\end{equation*}
Originally, $y$ has three freely selectable bits: $y_1$, $y_2$, and $y_3$. But the equation imposes one constraint on them. For example, if we choose $y_1$ and $y_2$ freely, then $y_3$ is determined by the equation:
\begin{equation*}
    y_3 = y_1 \oplus y_2
\end{equation*}
So only \textbf{two bits can be chosen independently}, and the third bit is determined by the first two. Hence, the dimension is governed by the number of independent bits:
\begin{equation*}
    \dim\left(a^{\perp}\right) = 3 - 1 = 2
\end{equation*}
In general, $y$ contains $n$ bits/degrees of freedom, but the condition of orthogonality:
\begin{equation*}
    a \cdot y = 0
\end{equation*}
Provides \textbf{one nontrivial linear constraint} because Simon guarantees that $a \neq 0$. Therefore, the dimension of the orthogonal component is:
\begin{equation*}
    n \text{ degrees of freedom} - 1 \text{ constraint} = n - 1
\end{equation*}
In our previous example, the constraint $a \cdot y = 0$ forces $y_3$ to be determined by $y_1$ and $y_2$, so we lose one degree of freedom, and the dimension is:
\begin{equation}
    \dim\left(a^{\perp}\right) = n - 1
\end{equation}
This is the \textbf{dimension of the set of all possible measurements} for a given hidden string $a$. In other words, it means that every possible measurement $y$ can be described using $n-1$ \textbf{independent binary choices}. Because each independent choice can be either $0$ or $1$, the \textbf{total number of elements} in the set is:
\begin{equation}
    \left|a^{\perp}\right| = 2^{n-1}
\end{equation}
For example, if $n=3$, and $a = 111$, the set of all possible measurements is:
\begin{equation*}
    a^{\perp} = \{000, 011, 101, 110\}
\end{equation*}
And its dimension is:
\begin{equation*}
    \dim\left(a^{\perp}\right) = 3 - 1 = 2
\end{equation*}
But its size (number of possible measurements) is:
\begin{equation*}
    \left|a^{\perp}\right| = 2^{3-1} = 4
\end{equation*}

\highspace
\textcolor{Green3}{\faIcon{link} \textbf{If we don't know $a$, but Simon gives us $y$'s, can we recover $a$?}} So far we have been looking at the problem from the perspective:
\begin{equation*}
    a \text{ is fixed, and we are finding all possible } y \text{'s}
\end{equation*}
We found that the set of all possible measurements is:
\begin{equation*}
    a^{\perp} = \left\{
        y \in \left\{0,1\right\}^{n} \mid a \cdot y = 0
    \right\}
\end{equation*}
And we found that its dimension is:
\begin{equation*}
    \dim\left(a^{\perp}\right) = n - 1
\end{equation*}
Now, here, we \textbf{reverse the direction}. Suppose we run the circuit several times and obtain:
\begin{equation*}
    y^{(1)}, y^{(2)}, \ldots, y^{(m)}
\end{equation*}
Every one belongs to $a^{\perp}$, so they all satisfy:
\begin{equation*}
    a \cdot y^{(i)} = 0 \quad \text{for } i = 1, 2, \ldots, m
\end{equation*}
We know that the dimension of $a^{\perp}$ is:
\begin{equation*}
    \dim\left(a^{\perp}\right) = n - 1
\end{equation*}
That means there can be $n-1$ \textbf{independent directions} inside the space of valid measurements. So if Simon eventually gives us $n-1$ \textbf{linearly independent} measurements:
\begin{equation*}
    y^{(1)}, y^{(2)}, \ldots, y^{(n-1)}
\end{equation*}
They collectively describe the whole space of valid measurements $a^{\perp}$. And now comes the key reversal: \hl{\emph{if we know the entire space $a^{\perp}$, which is the vector i orthogonal to \textbf{all} vectors in that space?}} The answer is: \textbf{the hidden string $a$ itself}. In other words, if we know the entire space of valid measurements, we can recover the hidden string $a$.

\begin{examplebox}[: $n=3$ and $a = 111$]
    Suppose we have $n=3$, and the hidden string is unknown:
    \begin{equation*}
        a = a_1 a_2 a_3
    \end{equation*}
    \textbf{First run.} Suppose Simon measures;
    \begin{equation*}
        y^{(1)} = 110
    \end{equation*}
    We know from everything we proved earlier that every measurement must satisfy:
    \begin{equation*}
        a \cdot y = 0
    \end{equation*}
    Therefore, the first measurement gives us the equation:
    \begin{equation*}
        110 \cdot a = 0
    \end{equation*}
    Which gives:
    \begin{equation*}
        1 a_1 \oplus 1 a_2 \oplus 0 a_3 = a_1 \oplus a_2 = 0
    \end{equation*}
    From XOR properties, $a_1 \oplus a_2 = 0$ means that $a_1 = a_2$. But that's not enough to determine $a$ yet. We need more measurements.

    \highspace
    But suppose we would like to stop here and try to guess $a$. We know that $a_1 = a_2$:
    \begin{equation*}
        \begin{cases}
            a_1 = a_2 \\
            a_3 \text{ is free}
        \end{cases}
    \end{equation*}
    This means that $a$ can be either:
    \begin{equation*}
        a = 000, \quad a = 110, \quad a = 001, \quad a = 111
    \end{equation*}
    But Simon guarantees that $a \neq 0$, so we can eliminate $000$.
    So we are left with three possibilities:
    \begin{equation*}
        a = 110, \quad a = 001, \quad a = 111
    \end{equation*}
    At this point, we cannot determine $a$ uniquely. We need more measurements!

    \highspace
    \textbf{Second run.} Suppose we run Simon again measure:
    \begin{equation*}
        y^{(2)} = 011
    \end{equation*}
    This gives us the equation:
    \begin{equation*}
        011 \cdot a = 0
    \end{equation*}
    Which gives:
    \begin{equation*}
        0 a_1 \oplus 1 a_2 \oplus 1 a_3 = a_2 \oplus a_3 = 0
    \end{equation*}
    From XOR properties, $a_2 \oplus a_3 = 0$ means that $a_2 = a_3$. Now we have two equations:
    \begin{equation*}
        \begin{cases}
            a_1 \oplus a_2 = 0 \\
            a_2 \oplus a_3 = 0
        \end{cases}
        \Rightarrow
        \begin{cases}
            a_1 = a_2 \\
            a_2 = a_3
        \end{cases}
        \Rightarrow a_1 = a_2 = a_3
    \end{equation*}
    So the hidden string must be either $000$ or $111$. But Simon guarantees that $a \neq 0$, so the only solution is:
    \begin{equation*}
        a = 111
    \end{equation*}
    And indeed, if we check the dimension of the space of valid measurements, we find that it is $n-1 = 2$, so two independent measurements are enough to recover $a$.
\end{examplebox}